\begin{tikzpicture}[scale=2]
    % 定义参数
    \def\a{1.8} % 圆心 y 坐标
    \def\thetaAngle{40} % 半径与负 y 轴的夹角 (对应草图中的 \theta)
    
    % 坐标轴
    \draw[->, thick] (-2, 0) -- (2.5, 0) node[below] {$x$};
    \draw[->, thick] (0, -0.5) -- (0, 3.5) node[right] {$y$};
    \node[below left] at (0,0) {$O$};
    
    % 抛物线 \Gamma_1: y=x^2
    \draw[thick, name path=parabola] plot[domain=-1.5:1.6, samples=100] (\x, {\x*\x});
    
    % 圆心 A(0, a)
    \coordinate (A) at (0, \a);
    \fill (A) circle (1.2pt) node[left] {$A(0,a)$};
    
    % 圆 \Gamma_2 (半径为1)
    \draw[thick] (A) circle (1);
    
    % 切点 P: (\sin\theta, a - \cos\theta)
    \coordinate (P) at ({sin(\thetaAngle)}, {\a - cos(\thetaAngle)});
    
    % 计算切线的斜率和截距
    % 根据几何关系，半径 AP 的倾角为 -90+\thetaAngle，因此切线倾角为 \thetaAngle
    \def\slope{tan(\thetaAngle)}
    \def\yint{\a - 1/cos(\thetaAngle)}
    
    % 绘制切线 (红色)
    \draw[thick, black, name path=tangent] plot[domain=-1.3:1.8] (\x, {\slope*\x + \yint});
    
    % 半径 AP (红色)
    \draw[thick, black] (A) -- (P);
    
    % 辅助线：向下垂线和 P 点水平线
    \coordinate (Adown) at (0, {\a - 1.2});
    \draw[dashed] (A) -- (Adown);
    \coordinate (Pright) at ($(P) + (1, 0)$);
    \draw[dashed] (P) -- (Pright);
    
    % 使用 intersections 求切线与抛物线的交点 M, N
    \path [name intersections={of=parabola and tangent, by={M,N}}];
    
    % 绘制交点及标签 (红色)
    \fill[black] (P) circle (1.2pt) node[below] {$M$};
    \fill[black] (M) circle (1.2pt) node[left] {$P$};
    \fill[black] (N) circle (1.2pt) node[right] {$Q$};
    
    % 标记角度 \theta
    % A 处的 \theta
    \pic [draw, black, "$\theta$", angle radius=0.4cm, angle eccentricity=1.5] {angle = Adown--A--P};
    
    % P 处的 \theta (切线与水平线的夹角)
    \coordinate (Pline) at ($(P) + (1, {\slope})$);
    \pic [draw, black, "$\theta$", angle radius=0.5cm, angle eccentricity=1.4] {angle = Pright--P--Pline};
    
    % 标记切点 P 处的直角
    \coordinate (Ptang) at ($(P) + (-0.5, {-0.5*\slope})$);
    \pic [draw, black, angle radius=0.2cm] {right angle = A--P--Ptang};
    
\end{tikzpicture}